% from ega0-19-fr.tex
\oldpage[0\textsubscript{IV-1}]{85}

\begin{proof}[Suite de la démonstration du lemme~{\rm(19.3.10.1)}]
Appliquant ce lemme à la situation de (19.3.10.1), la commutativité du
diagramme
\[
\xymatrix{
  B \ar[r]^{v'} \ar[d]_u & E'=E/\mathfrak L \ar[d] \\
  E/\mathfrak R \ar[r]_{f'} & E'/\mathfrak R'
}
\]
montre l'existence d'un homomorphisme
$w:B\to E/(\mathfrak R\cap\mathfrak L)$ tel que $v'$ et $u$ se
factorisent en
$B\xrightarrow{w}E/(\mathfrak R\cap\mathfrak L)\to E/\mathfrak L$
et
$B\xrightarrow{w}E/(\mathfrak R\cap\mathfrak L)\to E/\mathfrak R$
respectivement~; comme le noyau de $w$ contient l'intersection de ceux de
$v'$ et de $u$, il est ouvert dans $B$ et $w$ est continu. Enfin, il est clair
que $\mathfrak R\cap\mathfrak L$ est nilpotent, donc $w$ se factorise en
$B\xrightarrow{v}E\to E/(\mathfrak R\cap\mathfrak L)$, où $v$ est
un $A$-homomorphisme continu répondant à la question.
\end{proof}

\begin{proof}[Suite de la démonstration de la proposition~{\rm(19.3.10)}]
L'existence de $v_{n+1}$ résulte du lemme (19.3.10.1) appliqué en remplaçant
$E$, $E'$, $\mathfrak R$, $\mathfrak R'$ par
$C/\mathfrak Q_{n+1}$, $C/\mathfrak Q_n$,
$(\mathfrak J+\mathfrak Q_{n+1})/\mathfrak Q_{n+1}$ et
$(\mathfrak J+\mathfrak Q_n)/\mathfrak Q_n$ respectivement, $u$, $u'$
et $v'$ par $u_{n+1}$, $u_n$ et $v_n$ respectivement. Notons que, comme $C$
est séparé et complet, il est égal à $\varprojlim_n(C/\mathfrak Q_n)$, et
$v=\varprojlim_n v_n$ est un $A$-homomorphisme continu de $B$ dans $C$.
Comme $\mathfrak J$ est fermé dans $C$, on a
$\mathfrak J=\bigcap_n(\mathfrak J+\mathfrak Q_n)$~; comme
$C/\mathfrak J$ est métrisable et complet et que les
$(\mathfrak J+\mathfrak Q_n)/\mathfrak J$ forment un système fondamental
de voisinages de $0$ dans $C/\mathfrak J$, on a
\[
  C/\mathfrak J=\varprojlim_n C/(\mathfrak J+\mathfrak Q_n)
\]
et $\varprojlim_n\varphi_n$ est l'application canonique
$\varphi:C\to C/\mathfrak J$. Comme $\varprojlim_n u_n=u$, on a bien
$u=v\circ\varphi$, C.Q.F.D.
\end{proof}

\begin{corollary}[19.3.11]
\label{0.19.3.11-fr}
Soient $A$ un anneau topologique, $B$ une $A$-algèbre formellement lisse,
$C$ un anneau local noethérien complet, $\mathfrak J$ un idéal distinct de
$C$, $\varphi:A\to C$ un homomorphisme continu, faisant de $C$ une
$A$-algèbre topologique. Alors tout $A$-homomorphisme continu
$u_0:B\to C_0=C/\mathfrak J$ se factorise en
$B\xrightarrow{u}C\to C/\mathfrak J$, où $u$ est un $A$-homomorphisme
continu.
\end{corollary}

Toutes les conditions de (19.3.10) sont en effet remplies
($\mathbf 0_{\rm I}$, 7.3.5).

\begin{proposition}[19.3.12]
\label{0.19.3.12-fr}
Soient $A$ un anneau topologique, $B$ une $A$-algèbre formellement lisse,
$C$ une $A$-algèbre topologique, $\mathfrak J$ un idéal de $C$, vérifiant les
conditions suivantes~:
\begin{enumerate}
  \item[$1^\circ$] Il existe un système fondamental d'idéaux ouverts
  $\mathfrak Q_\lambda$ de $C$, tels que les
  $C_\lambda=C/\mathfrak Q_\lambda$ soient des anneaux artiniens et que
  l'homomorphisme canonique $C\to\varprojlim_\lambda C_\lambda$ soit un
  isomorphisme d'anneaux topologiques.
  \item[$2^\circ$] L'idéal $\mathfrak J$ est fermé dans $C$ et
  topologiquement nilpotent.
  \item[$3^\circ$] Le carré de tout idéal ouvert de $B$ est ouvert.
\end{enumerate}
Dans ces conditions, tout $A$-homomorphisme continu
$u:B\to C/\mathfrak J$ se factorise en
$B\xrightarrow{v}C\to C/\mathfrak J$, où $v$ est un $A$-homomorphisme
continu.
\end{proposition}

\begin{proof}
Soit $\mathfrak J_\lambda=(\mathfrak J+\mathfrak Q_\lambda)/\mathfrak Q_\lambda$
l'image canonique de $\mathfrak J$ dans $C_\lambda$, qui est un idéal
nilpotent, et soit $u_\lambda$ l'homomorphisme composé
$B\xrightarrow{u}C/\mathfrak J\to C_\lambda/\mathfrak J_\lambda
=C/(\mathfrak J+\mathfrak Q_\lambda)$. Montrons que tout
$A$-homomorphisme $w:B\to C$ tel que $u$ se factorise en
$B\xrightarrow{w}C\to C/\mathfrak J$ est nécessairement continu~; en
effet, pour tout $\lambda$, il y a un idéal ouvert $\mathfrak R$ de $B$
tel que $u(\mathfrak R)\subset
(\mathfrak J+\mathfrak Q_\lambda)/\mathfrak J$, d'où
$w(\mathfrak R)\subset\mathfrak J+\mathfrak Q_\lambda$, donc il y a
une puissance $\mathfrak R^{2^m}$ de $\mathfrak R$ telle que
$w(\mathfrak R^{2^m})\subset\mathfrak Q_\lambda$, ce qui établit
notre assertion puisque $\mathfrak R^{2^m}$ est ouvert dans $B$. On peut
donc se borner à trouver un $A$-homomorphisme $w$ ayant la propriété précédente
sans s'inquiéter de ses propriétés de continuité. Or, l'ensemble
$\operatorname{Hom}(B,C)$ de tous les $A$-homomorphismes d'algèbres de $B$
dans $C$ est égal à
$\varprojlim_\lambda\operatorname{Hom}(B,C_\lambda)$, et ce dernier est
\emph{fermé} dans le $C_\lambda$-module $C_\lambda^B$, muni de la topologie
produit, pour laquelle il est \emph{linéairement compact} puisque
  $C_\lambda$ est artinien. Pour tout $\lambda$, soit $W_\lambda$

